\documentclass[12pt,letterpaper]{article}
\usepackage{graphicx,textcomp}
\usepackage{natbib}
\usepackage{setspace}
\usepackage{fullpage}
\usepackage{color}
\usepackage[reqno]{amsmath}
\usepackage{amsthm}
\usepackage{fancyvrb}
\usepackage{amssymb,enumerate}
\usepackage[all]{xy}
\usepackage{endnotes}
\usepackage{lscape}
\newtheorem{com}{Comment}
\usepackage{float}
\usepackage{hyperref}
\newtheorem{lem} {Lemma}
\newtheorem{prop}{Proposition}
\newtheorem{thm}{Theorem}
\newtheorem{defn}{Definition}
\newtheorem{cor}{Corollary}
\newtheorem{obs}{Observation}
\usepackage[compact]{titlesec}
\usepackage{dcolumn}
\usepackage{tikz}
\usetikzlibrary{arrows}
\usepackage{multirow}
\usepackage{xcolor}
\newcolumntype{.}{D{.}{.}{-1}}
\newcolumntype{d}[1]{D{.}{.}{#1}}
\definecolor{light-gray}{gray}{0.65}
\usepackage{url}
\usepackage{listings}
\usepackage{color}

\definecolor{codegreen}{rgb}{0,0.6,0}
\definecolor{codegray}{rgb}{0.5,0.5,0.5}
\definecolor{codepurple}{rgb}{0.58,0,0.82}
\definecolor{backcolour}{rgb}{0.95,0.95,0.92}

\lstdefinestyle{mystyle}{
	backgroundcolor=\color{backcolour},
	commentstyle=\color{codegreen},
	keywordstyle=\color{magenta},
	numberstyle=\tiny\color{codegray},
	stringstyle=\color{codepurple},
	basicstyle=\footnotesize,
	breakatwhitespace=false,
	breaklines=true,
	captionpos=b,
	keepspaces=true,
	numbers=left,
	numbersep=5pt,
	showspaces=false,
	showstringspaces=false,
	showtabs=false,
	tabsize=2
}
\lstset{style=mystyle}
\newcommand{\Sref}[1]{Section~\ref{#1}}
\newtheorem{hyp}{Hypothesis}

\title{Problem Set 3}
\date{Due: March 25, 2026}
\author{Applied Stats II}


\begin{document}
	\maketitle
	\section*{Instructions}
	\begin{itemize}
	\item Please show your work! You may lose points by simply writing in the answer. If the problem requires you to execute commands in \texttt{R}, please include the code you used to get your answers. Please also include the \texttt{.R} file that contains your code. If you are not sure if work needs to be shown for a particular problem, please ask.
\item Your homework should be submitted electronically on GitHub in \texttt{.pdf} form.
\item This problem set is due before 23:59 on Wednesday March 25, 2026. No late assignments will be accepted.
	\end{itemize}

	\vspace{.25cm}
\section*{Question 1}
\vspace{.25cm}
\noindent We are interested in how governments' management of public resources impacts economic prosperity. Our data come from \href{https://www.researchgate.net/profile/Adam_Przeworski/publication/240357392_Classifying_Political_Regimes/links/0deec532194849aefa000000/Classifying-Political-Regimes.pdf}{Alvarez, Cheibub, Limongi, and Przeworski (1996)} and is labelled \texttt{gdpChange.csv} on GitHub. The dataset covers 135 countries observed between 1950 or the year of independence or the first year forwhich data on economic growth are available ("entry year"), and 1990 or the last year for which data on economic growth are available ("exit year"). The unit of analysis is a particular country during a particular year, for a total $>$ 3,500 observations.

\begin{itemize}
	\item
	Response variable:
	\begin{itemize}
		\item \texttt{GDPWdiff}: Difference in GDP between year $t$ and $t-1$. Possible categories include: "positive", "negative", or "no change"
	\end{itemize}
	\item
	Explanatory variables:
	\begin{itemize}
		\item
		\texttt{REG}: 1=Democracy; 0=Non-Democracy
		\item
		\texttt{OIL}: 1=if the average ratio of fuel exports to total exports in 1984-86 exceeded 50\%; 0= otherwise
	\end{itemize}

\end{itemize}
\newpage
\noindent Please answer the following questions:

\begin{enumerate}
	\item Construct and interpret an unordered multinomial logit with \texttt{GDPWdiff} as the output and "no change" as the reference category, including the estimated cutoff points and coefficients.

\vspace{.25cm}

First, we prepare the data by converting \texttt{GDPWdiff} to a factor with ``no change'' as the reference level, and converting \texttt{REG} and \texttt{OIL} to labelled factors:

\lstinputlisting[language=R, firstline=49, lastline=77]{PS3_SN.R}
\vspace{.5cm}


% Table created by stargazer v.5.2.3 by Marek Hlavac, Social Policy Institute. E-mail: marek.hlavac at gmail.com
% Date and time: Sun, Mar 22, 2026 - 21:07:41
\begin{table}[!htbp] \centering 
  \caption{} 
  \label{} 
\begin{tabular}{@{\extracolsep{5pt}}lcc} 
\\[-1.8ex]\hline 
\hline \\[-1.8ex] 
 & \multicolumn{2}{c}{\textit{Dependent variable:}} \\ 
\cline{2-3} 
\\[-1.8ex] & negative & positive \\ 
\\[-1.8ex] & (1) & (2)\\ 
\hline \\[-1.8ex] 
 REGDemocracy & 1.379$^{*}$ & 1.769$^{**}$ \\ 
  & (0.769) & (0.767) \\ 
  & & \\ 
 OILOil Exporter & 4.784 & 4.576 \\ 
  & (6.885) & (6.885) \\ 
  & & \\ 
 Constant & 3.805$^{***}$ & 4.534$^{***}$ \\ 
  & (0.271) & (0.269) \\ 
  & & \\ 
\hline \\[-1.8ex] 
Akaike Inf. Crit. & 4,690.770 & 4,690.770 \\ 
\hline 
\hline \\[-1.8ex] 
\textit{Note:}  & \multicolumn{2}{r}{$^{*}$p$<$0.1; $^{**}$p$<$0.05; $^{***}$p$<$0.01} \\ 
\end{tabular} 
\end{table}  


\vspace{.5cm}
We also compute the relative risk ratios (RRRs) by exponentiating the coefficients:


% Table created by stargazer v.5.2.3 by Marek Hlavac, Social Policy Institute. E-mail: marek.hlavac at gmail.com
% Date and time: Sun, Mar 22, 2026 - 21:07:41
\begin{table}[!htbp] \centering 
  \caption{Relative Risk Ratios (Unordered Multinomial Logit)} 
  \label{} 
\begin{tabular}{@{\extracolsep{5pt}} ccccc} 
\\[-1.8ex]\hline 
\hline \\[-1.8ex] 
 & (Intercept) & REGDemocracy & OILOil Exporter & Category \\ 
\hline \\[-1.8ex] 
negative & $44.942$ & $3.972$ & $119.578$ & negative \\ 
positive & $93.108$ & $5.865$ & $97.156$ & positive \\ 
\hline \\[-1.8ex] 
\end{tabular} 
\end{table}  


\subsection*{Interpretation}

The unordered multinomial logit produces two separate equations, one for each non-reference category compared to ``no change'':

\textbf{Intercepts (not cutoff points):} Unlike the ordered model, the unordered multinomial logit has no latent scale, so it does not estimate cutoff points. Instead, each equation has its own \textit{intercept} giving the baseline log-odds of that category versus the reference (``no change'') when all predictors are at their reference level (non-democracy, not an oil exporter).

\textbf{REG (Democracy):} The coefficient for Democracy in the ``negative'' equation gives the change in log-odds of experiencing negative GDP growth (vs.\ no change) when a country is a democracy rather than a non-democracy. Similarly, the coefficient in the ``positive'' equation gives the change in log-odds of experiencing positive GDP growth (vs.\ no change). The RRR for Democracy is 3.97 in the ``negative'' row and 5.87 in the ``positive'' row. This means that, relative to no change, democracies have a relative risk of negative GDP change that is 3.97 times higher than non-democracies, and a relative risk of positive GDP change that is 5.87 times higher. Since the ``positive'' RRR is larger, democracies are shifted more toward positive GDP outcomes relative to the ``no change'' baseline.

\textbf{OIL (Oil Exporter):} The RRR for Oil Exporter is 119.58 in the ``negative'' row and 97.16 in the ``positive'' row. Both are very large, meaning oil exporters are far more likely to experience \textit{either} negative or positive GDP change compared to no change. The ``negative'' RRR being slightly larger suggests oil exporters are marginally more tilted toward negative GDP change relative to positive. However, these very large RRRs reflect the fact that ``no change'' is an extremely rare category in these data (only 16 of 3,721 observations), making the reference group very small and the ratios unstable.

\newpage

	\item Construct and interpret an ordered multinomial logit with \texttt{GDPWdiff} as the outcome variable, including the estimated cutoff points and coefficients.

\vspace{.25cm}

For the ordered model, we create an ordered factor with the natural ranking negative $<$ no change $<$ positive:

\lstinputlisting[language=R, firstline=79, lastline=99]{PS3_SN.R}
\vspace{.5cm}


% Table created by stargazer v.5.2.3 by Marek Hlavac, Social Policy Institute. E-mail: marek.hlavac at gmail.com
% Date and time: Sun, Mar 22, 2026 - 21:07:41
\begin{table}[!htbp] \centering 
  \caption{} 
  \label{} 
\begin{tabular}{@{\extracolsep{5pt}}lc} 
\\[-1.8ex]\hline 
\hline \\[-1.8ex] 
 & \multicolumn{1}{c}{\textit{Dependent variable:}} \\ 
\cline{2-2} 
\\[-1.8ex] & GDPWdiff\_ord \\ 
\hline \\[-1.8ex] 
 REGDemocracy & 0.398$^{***}$ \\ 
  & (0.075) \\ 
  & \\ 
 OILOil Exporter & $-$0.199$^{*}$ \\ 
  & (0.116) \\ 
  & \\ 
\hline \\[-1.8ex] 
Observations & 3,721 \\ 
\hline 
\hline \\[-1.8ex] 
\textit{Note:}  & \multicolumn{1}{r}{$^{*}$p$<$0.1; $^{**}$p$<$0.05; $^{***}$p$<$0.01} \\ 
\end{tabular} 
\end{table}  


\subsection*{Interpretation}

The ordered logit (proportional odds) model assumes a latent continuous variable underlying the three ordered categories, with two cutoff points dividing the latent scale:

\textbf{Cutoff points (thresholds):}
\begin{itemize}
	\item \texttt{negative|no change}: the threshold on the latent scale separating ``negative'' from ``no change.'' Observations with latent values below this threshold are predicted as ``negative.''
	\item \texttt{no change|positive}: the threshold separating ``no change'' from ``positive.'' Observations with latent values above this threshold are predicted as ``positive.''
\end{itemize}

\textbf{REG (Democracy):} A positive coefficient means that being a democracy shifts the latent variable upward, increasing the probability of being in a higher category (toward ``positive'' GDP change) and decreasing the probability of being in a lower category (``negative''). The proportional odds assumption means this single coefficient applies equally at both thresholds. The odds ratio $e^{\beta_{\text{REG}}}$ gives the multiplicative change in the odds of being in a higher category for democracies versus non-democracies.

\textbf{OIL (Oil Exporter):} Similarly, the OIL coefficient captures the effect of being an oil exporter on the latent GDP change variable.

\textbf{Key difference from the unordered model:} The ordered model imposes the proportional odds assumption --- a single coefficient per predictor rather than separate coefficients for each category comparison. This is more parsimonious but assumes the effect is uniform across the ordered thresholds.


\end{enumerate}

\newpage

\section*{Question 2}
\vspace{.25cm}

\noindent Consider the data set \texttt{MexicoMuniData.csv}, which includes municipal-level information from Mexico. The outcome of interest is the number of times the winning PAN presidential candidate in 2006 (\texttt{PAN.visits.06}) visited a district leading up to the 2009 federal elections, which is a count. Our main predictor of interest is whether the district was highly contested, or whether it was not (the PAN or their opponents have electoral security) in the previous federal elections during 2000 (\texttt{competitive.district}), which is binary (1=close/swing district, 0="safe seat"). We also include \texttt{marginality.06} (a measure of poverty) and \texttt{PAN.governor.06} (a dummy for whether the state has a PAN-affiliated governor) as additional control variables.

\begin{enumerate}
	\item [(a)]
	Run a Poisson regression because the outcome is a count variable. Is there evidence that PAN presidential candidates visit swing districts more? Provide a test statistic and p-value.

\vspace{.25cm}

\lstinputlisting[language=R, firstline=113, lastline=129]{PS3_SN.R}
\vspace{.5cm}


% Table created by stargazer v.5.2.3 by Marek Hlavac, Social Policy Institute. E-mail: marek.hlavac at gmail.com
% Date and time: Sun, Mar 22, 2026 - 21:07:41
\begin{table}[!htbp] \centering 
  \caption{} 
  \label{} 
\begin{tabular}{@{\extracolsep{5pt}}lc} 
\\[-1.8ex]\hline 
\hline \\[-1.8ex] 
 & \multicolumn{1}{c}{\textit{Dependent variable:}} \\ 
\cline{2-2} 
\\[-1.8ex] & PAN.visits.06 \\ 
\hline \\[-1.8ex] 
 competitive.district & $-$0.081 \\ 
  & (0.171) \\ 
  & \\ 
 marginality.06 & $-$2.080$^{***}$ \\ 
  & (0.117) \\ 
  & \\ 
 PAN.governor.06 & $-$0.312$^{*}$ \\ 
  & (0.167) \\ 
  & \\ 
 Constant & $-$3.810$^{***}$ \\ 
  & (0.222) \\ 
  & \\ 
\hline \\[-1.8ex] 
Observations & 2,407 \\ 
Log Likelihood & $-$645.606 \\ 
Akaike Inf. Crit. & 1,299.213 \\ 
\hline 
\hline \\[-1.8ex] 
\textit{Note:}  & \multicolumn{1}{r}{$^{*}$p$<$0.1; $^{**}$p$<$0.05; $^{***}$p$<$0.01} \\ 
\end{tabular} 
\end{table}  

\vspace{.5cm}

The Poisson model is:

$$\log(\mu_i) = \beta_0 + \beta_1 \cdot \texttt{competitive.district}_i + \beta_2 \cdot \texttt{marginality.06}_i + \beta_3 \cdot \texttt{PAN.governor.06}_i$$

The coefficient on \texttt{competitive.district} tests whether swing districts receive more visits:


% Table created by stargazer v.5.2.3 by Marek Hlavac, Social Policy Institute. E-mail: marek.hlavac at gmail.com
% Date and time: Sun, Mar 22, 2026 - 21:07:41
\begin{table}[!htbp] \centering 
  \caption{Test for Competitive District Effect} 
  \label{} 
\begin{tabular}{@{\extracolsep{5pt}} ccc} 
\\[-1.8ex]\hline 
\hline \\[-1.8ex] 
 & Statistic & Value \\ 
\hline \\[-1.8ex] 
1 & Coefficient & $$-$0.081$ \\ 
2 & z-statistic & $$-$0.477$ \\ 
3 & p-value & $0.634$ \\ 
4 & IRR & $0.922$ \\ 
\hline \\[-1.8ex] 
\end{tabular} 
\end{table}  

\vspace{.5cm}

If $\hat{\beta}_1 > 0$ and $p < 0.05$, we have evidence that PAN presidential candidates visit competitive districts more frequently, holding poverty level and governor affiliation constant. The incidence rate ratio $e^{\hat{\beta}_1}$ gives the multiplicative increase in the expected visit count for competitive districts relative to safe seats.

	\item [(b)]
	Interpret the \texttt{marginality.06} and \texttt{PAN.governor.06} coefficients.

\vspace{.25cm}

\lstinputlisting[language=R, firstline=131, lastline=137]{PS3_SN.R}
\vspace{.5cm}


% Table created by stargazer v.5.2.3 by Marek Hlavac, Social Policy Institute. E-mail: marek.hlavac at gmail.com
% Date and time: Sun, Mar 22, 2026 - 21:07:41
\begin{table}[!htbp] \centering 
  \caption{Incidence Rate Ratios (Poisson)} 
  \label{} 
\begin{tabular}{@{\extracolsep{5pt}} cc} 
\\[-1.8ex]\hline 
\hline \\[-1.8ex] 
 & IRR \\ 
\hline \\[-1.8ex] 
(Intercept) & $0.022$ \\ 
competitive.district & $0.922$ \\ 
marginality.06 & $0.125$ \\ 
PAN.governor.06 & $0.732$ \\ 
\hline \\[-1.8ex] 
\end{tabular} 
\end{table}  

\vspace{.5cm}

\textbf{marginality.06:} For a one-unit increase in the marginality index (i.e., moving toward greater poverty), the expected number of PAN candidate visits is multiplied by the IRR for marginality, holding other variables constant. If the coefficient is negative (IRR $<$ 1), poorer districts receive fewer visits; if positive (IRR $>$ 1), they receive more.

\textbf{PAN.governor.06:} Districts in states with a PAN-affiliated governor have an expected visit count that is multiplied by the IRR for PAN governor relative to districts without one, holding competitiveness and marginality constant. An IRR $>$ 1 indicates that PAN candidates visit districts with a co-partisan governor more frequently, which may reflect strategic coordination or stronger party infrastructure.

	\item [(c)]
	Provide the estimated mean number of visits from the winning PAN presidential candidate for a hypothetical district that was competitive (\texttt{competitive.district}=1), had an average poverty level (\texttt{marginality.06} = 0), and a PAN governor (\texttt{PAN.governor.06}=1).

\vspace{.25cm}

\lstinputlisting[language=R, firstline=139, lastline=151]{PS3_SN.R}
\vspace{.5cm}

The predicted count is calculated as:

$$\hat{\mu} = e^{\hat{\beta}_0 + \hat{\beta}_1(1) + \hat{\beta}_2(0) + \hat{\beta}_3(1)} = e^{\hat{\beta}_0 + \hat{\beta}_1 + \hat{\beta}_3}$$

Since \texttt{marginality.06} = 0, its term drops out:


% Table created by stargazer v.5.2.3 by Marek Hlavac, Social Policy Institute. E-mail: marek.hlavac at gmail.com
% Date and time: Sun, Mar 22, 2026 - 21:07:41
\begin{table}[!htbp] \centering 
  \caption{Predicted Number of PAN Visits} 
  \label{} 
\begin{tabular}{@{\extracolsep{5pt}} ccc} 
\\[-1.8ex]\hline 
\hline \\[-1.8ex] 
 & Scenario & Predicted.Visits \\ 
\hline \\[-1.8ex] 
1 & competitive=1, marginality=0, PAN.governor=1 & $0.015$ \\ 
\hline \\[-1.8ex] 
\end{tabular} 
\end{table}  


\end{enumerate}

\section*{AI Usage Disclosure}
Claude Code (Anthropic, Claude Opus 4.6) was used to assist with R code development, \LaTeX{} formatting, and interpretation drafting for this problem set. Models used: Claude Opus 4.6 (1M context). All output was reviewed, verified, and edited by the author.

\end{document}
